% SOURCE_STATEMENT: "For the nontransient oscillatory sedimentation of a sphere at low Reynolds number in a wormlike micellar fluid, what constitutive equation or constitutive model gives a mathematical understanding of the observed oscillations?"
% Preparation and final self-audit: CONTRIBUTING.md prompts C and D; source/audit artifacts retained in the local physics_statements_2026_10_11 batch.
% !TeX program = lualatex
\documentclass[11pt]{article}
\usepackage[margin=25mm]{geometry}
\usepackage{fontspec}
\setmainfont{Latin Modern Roman}
\usepackage{amsmath,amssymb,mathtools}
\usepackage{unicode-math}
\setmathfont{Latin Modern Math}
\usepackage{xurl,hyperref}
\hypersetup{colorlinks=true,urlcolor=blue}
\setcounter{secnumdepth}{0}
\setlength{\parindent}{0pt}
\setlength{\parskip}{5pt}
\setlength{\emergencystretch}{3em}
\newcommand{\sref}[2]{\hyperlink{source:#1:#2}{[S#2]}}
\newcommand{\cataloguescope}[1]{\paragraph{Recorded target.}#1}
\newcommand{\R}{\mathbb R}
\newcommand{\C}{\mathbb C}
\newcommand{\N}{\mathbb N}
\newcommand{\Z}{\mathbb Z}
\newcommand{\Q}{\mathbb Q}
\newcommand{\F}{\mathbb F}
\newcommand{\eps}{\varepsilon}
\begin{document}
\section{Constitutive Models for Oscillatory Micellar Sedimentation}\label{statement}

\subsection{Definitions and mathematical statement}

\paragraph{Shared notation.}$\mathbb N=\{1,2,\ldots\}$; the standard number fields have their usual meanings. All additional conventions are stated locally.


\paragraph{Definitions and conventions.} The Reynolds number $\mathrm{Re}=\rho UL/\mu$ compares inertial and viscous effects. A wormlike micellar fluid contains long self-assembled surfactant aggregates whose breaking and recombination can affect viscoelastic stresses. A constitutive equation relates those stresses to deformation and internal structural variables. The observed phenomenon is sustained, rather than transient, oscillation of the settling velocity of a freely falling sphere at low Reynolds number.

\paragraph{Source scope and qualifications.} The target is the observed micellar sedimentation regime. The source gives no unique constitutive family, parameter identification, or mathematical oscillation criterion; importing a specific polymer model as the definition of success would narrow the question.

\paragraph{Catalogue mathematical target.}
\paragraph{Catalogue wording retained for traceability.}
For the nontransient oscillatory sedimentation of a sphere at low Reynolds number in a wormlike micellar fluid, what constitutive equation or constitutive model gives a mathematical understanding of the observed oscillations?

\paragraph{Saved literature context.} The supplied catalogue's assessment dated 2026-08-21 records: The report explicitly identifies the constitutive-equation-level understanding of oscillating sedimentation in complex fluids as open; later work studies particular fluids and models rather than settling a general constitutive law. The assessment is retained as dated provenance; it is not a new comprehensive literature or open-status certification.
\subsection{Short English statement}
Find a constitutive model that explains sustained oscillation of a sphere settling through a wormlike micellar fluid.
\subsection{Sources}
\begin{itemize}
\item[\hypertarget{source:statement:1}{S1}] ULAM.AI and the UnsolvedMath Contributors, UnsolvedMath, catalogue entry. Catalogue curation under CC BY 4.0; definitions and English exposition added locally.
\url{https://huggingface.co/datasets/ulamai/UnsolvedMath}.
\item[\hypertarget{source:statement:2}{S2}] Mini-Workshop: Theory and Numerics of Fluid-Solid Interaction (2007), Oberwolfach Reports, DOI 10.4171/owr/2007/53, pp. 3131--3180; official report source and mathematical context.
\url{https://doi.org/10.4171/owr/2007/53}.
\item[\hypertarget{source:statement:3}{S3}] Official EMS Press PDF of the same Oberwolfach report; PDF pages 11--12.
\url{https://ems.press/content/serial-article-files/46139}.
\end{itemize}
\label{end:statement}
\end{document}
